% vim: spl=pt
\begin{exercício}{Descontinuidades da família espectral}{ex38}
  Seja \(T\in\bounded(\hilbert)\) um operador auto-adjunto e \(\family{E_\lambda^T}{\lambda \in \mathbb{R}}\) sua família espectral. Mostre que \(\lambda \mapsto E^T_{(-\infty, \lambda]}\) tem uma descontinuidade em \(\lambda_0\) se e somente se \(\lambda_0 \in \sigma_p(T)\) e
  \begin{equation*}
    \ker(T - \lambda_0\unity) = \left(E^T_{(-\infty, \lambda_0]} - \lim_{\epsilon\to 0^{-}}{E_{(-\infty, \lambda_0 + \epsilon]}}\right)\hilbert.
  \end{equation*}
\end{exercício}
\begin{proof}[Resolução]
  Vamos primeiro provar o
  \begin{lemma}{Projetor espectral e núcleo da parte positiva de um operador auto-adjunto}{kernel_spectral_family}
    Para um operador auto-adjunto \(A \in \bounded(\hilbert)\) denotamos \(A_+ = \frac12 \abs{A} + \frac12A \in \bounded(\hilbert)\) como sua parte positiva. Sendo \(V \subset \hilbert\) um subespaço fechado, denotamos \(P_V : \hilbert \to \hilbert\) como o projetor ortogonal sobre \(V.\) Com essas notações, temos para um operador auto-adjunto \(A \in \bounded(\hilbert)\) que
    \begin{equation*}
      E^A_{(-\infty, \lambda_0]} = P_{\ker{(A - \lambda_0\unity)_+}}
    \end{equation*}
    para todo \(\lambda_0 \in \mathbb{R}.\)
  \end{lemma}
  \begin{proof}
    Consideremos a função limitada
    \begin{align*}
      f : \sigma(S) &\to \mathbb{R}_+\\
                \xi &\mapsto \frac12\abs{\xi} + \frac12\xi,
    \end{align*}
    onde \(S \in \bounded(\hilbert)\). Temos então
    \begin{equation*}
      S_+ = \int_{\sigma(S)} \dli{E^S_\lambda} f(\lambda),
    \end{equation*}
    portanto
    \begin{align*}
      x \in \ker{S_+} &\iff \inner*{x}{\int_{\sigma(S)}\dli{E_\lambda^S} f(\lambda)x} = 0\\
                      &\iff \int_{\mathbb{R}} \dli{\mu^S_x(\lambda)} f(\lambda) = 0\\
                      &\iff \int_0^\infty \dli{\mu^S_x(\lambda)} f(\lambda) = 0.
    \end{align*}
    Como a medida espectral é positiva e \(f\) tem imagem em \(\mathbb{R}_+,\) segue que
    \begin{equation*}
      x \in \ker{S_+} \iff \mu_x(0,\infty) = 0 \iff \inner{x}{E_{(0,\infty)}^Sx} = 0 \iff x \in \ker{E_{(0,\infty)}^S} \iff x \in \ran{E^S_{(-\infty, 0]}},
    \end{equation*}
    já que \(E^S_{(-\infty, 0]} + E^S_{(0,\infty)} = E^S_{\mathbb{R}} = \unity.\) Assim, \(\ker{S_+} = \ran{E^S_{(-\infty, 0]}}\) e concluímos que \(E^S_{(-\infty, 0]} = P_{\ker{S_+}}.\)

    Para transladar o resultado por \(\lambda_0 \in \mathbb{R},\) notemos que para um boreliano \(B \subset \mathbb{R}\) temos
    \begin{align*}
      E^{A - \lambda_0\unity}_B &= \sup\setc{f(A - \lambda_0 \unity)}{f \in \mathcal{C}(\sigma(A - \lambda_0\unity) : 0 \leq f \nearrow \chi_{B}}\\
                                &= \sup\setc{g(A)}{g \in \mathcal{C}(\sigma(A) : 0\leq g \nearrow \chi_{B + \lambda_0}}\\
                                &= E^A_{B + \lambda_0},
    \end{align*}
    portanto \(P_{\ker(A - \lambda_0\unity)_+} = E^{A - \lambda_0\unity}_{(-\infty, 0]} = E^A_{(-\infty, \lambda_0]}.\)
  \end{proof}

  Notemos que para todo \(\lambda_0 \in \mathbb{R}\) e para todos borelianos \(A, B\) com \(A \subset B\) temos
  \begin{equation*}
    \left(E^T_{B} - E^T_{A}\right)^2 = E^T_{B} + E^T_{A} - E^T_{B}E^T_{A} - E^T_{A} E^T_{B} = E^T_{B} + E^T_{A} - 2E^T_{A \cap B} = E^T_{B} - E^T_{A},
  \end{equation*}
  já que \(A \cap B = A,\) portanto dado \(\epsilon < 0\) o operador
  \begin{equation*}
    P_\epsilon = E^T_{(-\infty, \lambda_0]} - E^T_{(-\infty, \lambda_0 + \epsilon]}
  \end{equation*}
  é um projetor ortogonal para todo \(\lambda_0 \in \mathbb{R}\) com \(E^T_{(-\infty, \lambda_0]}\hilbert \supset E^T_{(-\infty, \lambda_0 + \epsilon]}\hilbert.\)

  Sendo \(P = \lim_{\epsilon\to0^{-}}{P_\epsilon} = E^{T}_{\set{\lambda_0}}\), queremos mostrar que \(\ran{P} = \ker{(T - \lambda_0\unity)}.\) Seja \(x \in \ker(T - \lambda_0\unity),\) então
  \begin{align*}
    \norm{(T - \lambda_0\unity)_+x}^2 &= \frac14\norm{\sqrt{(T - \lambda_0\unity)^2}x}^2\\
                                      &= \frac14\inner{x}{(T - \lambda_0\unity)^2x}\\
                                      &= 0
  \end{align*}
  isto é, \(x \in \ran{E^T_{(-\infty,\lambda_0]}}.\) Notemos também que
  \begin{equation*}
    \left[T - (\lambda_0 + \epsilon)\unity\right]x = -\epsilon x \implies \left[T - (\lambda_0 - \epsilon)\unity\right]_+x = f(-\epsilon)x = \abs{\epsilon} x
  \end{equation*}
  portanto \(x \notin \ran{E^T_{(-\infty, \lambda_0 + \epsilon]}}\) para todo \(\epsilon < 0.\) Temos, portanto,
  \begin{equation*}
    \forall \epsilon < 0 , \forall x \in \ker(T - \lambda_0 \unity) : P_\epsilon x = E^T_{(-\infty, \lambda_0]}x - E^T_{(-\infty, \lambda_0 + \epsilon]}x = E^T_{(-\infty, \lambda_0]}x = x
  \end{equation*}
  logo \(\ker(T - \lambda_0\unity) \subset \ran{P_\epsilon}\) sempre que \(\epsilon < 0,\) logo no limite \(\epsilon \to 0^{-}\) temos \(\ker(T - \lambda_0\unity) \subset \ran{P}.\) Seja agora \(x \in \ran{P_\epsilon},\) então como \(T - \lambda_0\unity\) é auto-adjunto, segue que
  \begin{align*}
    \norm{(T - \lambda_0\unity)x}^2 &= \inner{x}{(T - \lambda_0\unity)^2x}\\
                                  &= \inner{x}{(T - \lambda_0\unity)^2P_\epsilon x}\\
                                  &= \abs*{\int_{\lambda_0 + \epsilon}^{\lambda_0} \dli{\mu_x^T(\lambda)} (\lambda - \lambda_0)^2}\\
                                  &\leq \mu_x^T((\lambda_0 + \epsilon, \lambda_0)) \sup_{\lambda \in (\lambda_0 + \epsilon, \lambda_0)}{\abs{\lambda - \lambda_0}^2}\\
                                  &= \inner{x}{E^T_{(\lambda_0 + \epsilon, \lambda_0)}x} \abs{\epsilon}^2
  \end{align*}
  para todo \(\epsilon < 0.\) Agora, como
  \begin{equation*}
    E^T_{(\lambda_0 + \epsilon, \lambda_0)} + E^T_{\set{\lambda_0}} = P_\epsilon,
  \end{equation*}
  temos
  \begin{equation*}
    0 \leq \inner{x}{E^T_{(\lambda_0 + \epsilon, \lambda_0)}x} = \norm{x}^2 - \mu_x^T(\set{\lambda_0}),
  \end{equation*}
  logo
  \begin{equation*}
    \norm{(T - \lambda_0\unity)x} \leq \abs{\epsilon}\sqrt{\norm{x}^2 - \mu_x^T(\set{\lambda_0})}
  \end{equation*}
  para todo \(\epsilon < 0,\) e concluímos que \(\norm{(T - \lambda_0\unity)x} = 0,\) isto é, \(\ran{P} \subset \ker{(T - \lambda_0\unity)}.\)

  No caso em que \(\lambda_0 \in \sigma_p(T),\) existe \(x \in \ker(T - \lambda_0 \unity) \setminus \set{0}\) tal que \((T - \lambda_0 \unity)x = \lambda_0x,\) e, portanto, \(Px = x.\) Caso \(\lambda_0 \notin \sigma_p(T),\) temos \(P = 0.\) Isso nos informa que a família espectral apresenta descontinuidade (na topologia forte) apenas no espectro pontual, já que o limite com \(\epsilon \to 0^{+}\) de \(P_\epsilon\) é o operador nulo, pois
  \begin{equation*}
    E_{(-\infty,\lambda_0+\epsilon]} - E_{(-\infty, \lambda_0]} = E_{(\lambda_0, \lambda_0+\epsilon]}
  \end{equation*}
  e \(\lambda_0 \notin (\lambda_0, \lambda_0+\epsilon]\) para todo \(\epsilon > 0.\)
\end{proof}
